% Revised copy: legend abbreviations, suffixes, and all three observed-frame theoretical curves are explained explicitly.
% Use inside a LaTeX figure environment after \includegraphics.
% Requires natbib (or another bibliography package supporting \citep)
% and amssymb (used by a cited paper title containing \gtrsim); amsmath is used by the model-parameter notation.
\caption{\label{fig:uvlf-z6-12p5}
Rest-frame ultraviolet luminosity functions from $z=6$ to $z=12.5$.
Panels \textbf{(a)}--\textbf{(g)} show the binned number density
$\phi$ as a function of absolute UV magnitude $M_{\rm UV}$ at
$z=6$, 7, 8, 9, 10, 11, and 12.5, respectively.
The primary measurements are emphasized with filled black symbols:
circles show the HST determinations of Bouwens et al. at $z=6$--8
\citep{Bouwens2021UVLF}, and squares show the JWST PRIMER
determinations of Donnan et al. at $z=9$--12.5
\citep{Donnan2024PRIMER}.
Colored open symbols show independent recent measurements from
Adams et al. \citep{Adams2024EPOCHS}, Atek et al.
\citep{Atek2024UNCOVER}, McLeod et al. \citep{McLeod2024UVLF},
Robertson et al. \citep{Robertson2024JADES}, Franco et al.
\citep{Franco2025COSMOSWeb}, Harikane et al.
\citep{Harikane2025Spectroscopic}, Rojas-Ruiz et al.
\citep{RojasRuiz2025BoRG}, Chemerynska et al.
\citep{Chemerynska2026GLIMPSE}, Curtis-Lake et al.
\citep{CurtisLake2026JADESDR4}, and Whitler et al.
\citep{Whitler2026JADES}.
Legend labels combine an abbreviated first-author surname with the
two-digit publication year; \mbox{Chem.}, \mbox{C.-Lake}, and
\mbox{R.-Ruiz} abbreviate Chemerynska, Curtis-Lake, and Rojas-Ruiz,
respectively. Additional suffixes distinguish independently analysed
samples: \mbox{C.-Lake+26 D} and \mbox{C.-Lake+26 M} denote the
Curtis-Lake Deep and Medium tiers; the 814, 115, and 150 suffixes for
Franco+25 denote the F814W-, F115W-, and F150W-dropout selections;
the 115 and 150 suffixes for Whitler+26 denote the F115W- and
F150W-dropout selections; and 1747 and 2426 for R.-Ruiz+25 identify
the corresponding JWST GO programme samples. These suffixes are
sample identifiers, not redshifts or sample sizes.
Horizontal and vertical bars reproduce the reported magnitude and
number-density uncertainties from each source; their statistical
definitions are therefore source dependent and have not been
homogenized. For a one-sided limit, the vertical error-bar arm in the
arrow direction is suppressed to avoid overplotting, while any
reported arm in the opposite direction is retained; this display-only
convention does not alter the tabulated source uncertainties. Thin
open arrows are drawn collinearly with the vertical error-bar direction
at the same $M_{\rm UV}$: upward arrows denote lower limits on $\phi$,
whereas downward arrows denote upper limits.
The green open hexagon in panel \textbf{(b)} is the Curtis-Lake et al.
Deep/JWST lower limit at $z_{\rm eff}=7.091$,
$M_{\rm UV}=-17.53^{+0.12}_{-0.05}$, and
$\phi>5.5\times10^{-4}\ {\rm mag}^{-1}\,{\rm Mpc}^{-3}$; it is one
constraint rather than several overlapping measurements.
The Curtis-Lake Deep and Medium tiers are retained as independent
samples. One Harikane et al. lower-limit row at $z=7$,
$M_{\rm UV}=-23.7$, and $\phi=2.3\times10^{-9}\
{\rm mag}^{-1}\,{\rm Mpc}^{-3}$ is not displayed because its central
anchor lies below the fixed vertical range. Four Chemerynska et al.
points centred at $M_{\rm UV}=-15.5$ remain visible, with only the
positive horizontal-error portions outside the plotting range
clipped. Three black curves compare the theoretical UV luminosity functions. Both PBH models assume instantaneous merger of halo PBHs at $z=20$, followed by accretion growth, with initial merged mass $M_h f_{\rm PBH}$. The PBH models use total-matter fluctuations, while the No PBH reference retains its previous non-neutrino-matter convention. The solid curve, labelled Fixed SF, uses fixed Planck-2018 cosmology and the updated fixed astrophysical parameters with the supplied eight-parameter MAP vector. The dashed Free SF curve uses the corrected 16-parameter MAP vector supplied in the parameter image, allowing the astrophysical parameters to differ from the solid case. The dotted No PBH curve uses the corrected no-PBH astrophysical parameter vector supplied in the parameter image while setting $M_{\rm PBH}=f_{\rm PBH}=0$ and explicitly omitting both PBH power-spectrum and PBH accretion-luminosity contributions. All three curves are mapped back to the published observational frame using the inverse dust-extinction, magnitude-bin-width, and Alcock--Paczynski distance and volume transformations used for the fitting data; the inverse dust mapping is active only at $z\leq8$.}
